% sections/02_background.tex

\section{Background \& Threat Model}
\label{sec:background}

\subsection{LLM Skill Marketplaces}

LLM skill marketplaces are application stores specialized for language
model extensions. A \emph{skill} is a self-contained module---typically a
\texttt{SKILL.md} manifest with tool definitions and helper code---that an
LLM agent invokes during conversation or automated tasks.
Skills may be first-party (platform-developed) or third-party.

Users install skills through a package-manager command
(\eg \texttt{/install author/skill-name}) and thereafter the skill
executes with \emph{ambient authority}: persistent access to the
agent's context window, conversation history, tools, and any connected
data sources (filesystem, email, calendar, credentials).

\paragraph{Trust signals.}
Platforms expose signals intended to help users evaluate skills:
\begin{enumerate}
    \item A free-text \textbf{description} written by the developer;
    \item Optional structured \textbf{metadata} (YAML declarations of capability);
    \item Version \textbf{changelogs};
    \item A platform-operated \textbf{security scan badge};
    \item Aggregate \textbf{download and star counts}.
\end{enumerate}
Our analysis suggests that signals 1--4 often understate actual capabilities, while signal 5 may amplify misplaced trust in risky skills.

\paragraph{Auto-update semantics.}
Both \clawhub{} and \skillssh{} propagate version updates automatically
to all installed users by default.
Users who consented to a skill at version $v_i$ are silently upgraded to
$v_{i+1}$ without re-review, re-consent, or changelog notification.
ClawHub's platform documentation confirms that updates propagate within
24~hours of publisher push.
Version pinning is unavailable to standard users; skill developers have
unilateral write access to all installed instances.
This is the key structural enabler of T4 and T6 attacks.

\subsection{Threat Model}

We model four principals:

\textbf{Users} install skills expecting them to behave as described.
They cannot audit code before installation and rely on platform trust signals.

\textbf{Benign (careless) developers} omit capability documentation
without malicious intent (T1, T4). Their omissions violate informed consent.

\textbf{Deceptive developers} consciously misstate capabilities to
increase installs or bypass platform review (T2, T3, T5).

\textbf{Malicious actors} (compromised accounts or supply-chain
attackers) deliberately inject harmful code into skills (T6).

\paragraph{Attack surface.}
The core vulnerability is \emph{consent at installation vs.\ capability
at runtime}. Users who consented to a ``calendar scheduler'' have not
consented to credential exfiltration---even if that capability is
in the installed code. This gap is structurally analogous to Android
permission creep~\cite{felt2011android} but operates at the description
level rather than the manifest level: Android \emph{requires} capability
declarations; skill platforms do not.

\paragraph{Attacker goal.}
A T6 attacker seeks to: (a) establish trust via a benign skill version,
accumulate installs, then (b) inject malicious capability via an update,
exploiting auto-update propagation, then (c) exfiltrate credentials/data
before the malicious version is detected and patched.
The expected harm is proportional to the install count and the exposure
window---both of which grow as the platform grows.

\paragraph{Why undisclosure is the attacker's rational strategy.}
A reviewer may ask why attackers do not use more sophisticated obfuscation
(runtime payload fetching, dynamic code generation, indirect API calls).
We observe that simple undisclosure is the \emph{risk-minimal} approach:
complexity increases the probability of detection by
pattern matching, sandbox execution, or anomaly scoring.
T6 evidence supports this: 32/34 confirmed T6 skills use
straightforward API-level patterns---\texttt{os.environ}~(CRED\_READ),
\texttt{requests.post}~(NET\_WRITE), \texttt{os.remove}~(FILE\_DELETE)---
without obfuscation. The marginally harder bar of tool-level
undisclosure is sufficient because platform scanners check
\emph{declared intent} (descriptions and metadata), not runtime capability.

\paragraph{Consent model scope.}
We do \emph{not} claim that users systematically read skill descriptions
or form complete capability expectations from them.
Rather, we note that descriptions are the \emph{only structured,
publish-time mechanism} for capability disclosure available to users across
both platforms we studied (YAML metadata fill rate: 0.0\%;
\S\ref{sec:static}).
When this mechanism fails systematically at scale, users and platform
operators have no alternative disclosure channel---regardless of
whether individual users attend to descriptions.
The threat is structural, not behavioral.

\subsection{LSCPS Capability Taxonomy}

We categorize skill capabilities into five top-level families,
abbreviated \textbf{\lscps{}}
(\underline{L}ocal-shell, \underline{S}ecret, \underline{C}ommand-and-control,
\underline{P}ersistent-write, \underline{S}ystemic-autonomous).
Each top-level family is subdivided into 14 fine-grained sub-categories
that serve as the atomic detection units (see Appendix~\ref{app:lscps}):

\begin{itemize}[topsep=2pt,itemsep=2pt]
\item[\textbf{L}] \textbf{Local Shell execution}:
  \texttt{subprocess}, \texttt{os.system}, \texttt{eval()}, bash invocation
\item[\textbf{S}] \textbf{Secret/Credential access}:
  API keys, OAuth tokens, passwords, private keys
\item[\textbf{C}] \textbf{Command-and-control / Network write}:
  \texttt{requests.post}, webhooks, SMTP exfiltration
\item[\textbf{P}] \textbf{Persistent file write}:
  \texttt{open(w)}, \texttt{shutil.rmtree}, \texttt{os.remove}
\item[\textbf{S}] \textbf{Systemic/Autonomous effects}:
  send email, create/modify files, background scheduling, multi-step no-confirm loops
\end{itemize}

Each sub-category is detected by a curated regex bank totaling 300+ patterns
(see Appendix~\ref{app:lscps}),
validated on the 40 T6 ground-truth version pairs
(72.5\% pair-level high-risk recall; 79.4\% at skill level,
\S\ref{sec:e2e}).

\paragraph{Methodological rationale: Why LSCPS, not Bandit/Semgrep?}
General-purpose security linters (Bandit, Semgrep) detect Python anti-patterns
(SQL injection, hardcoded secrets) but are designed for \emph{single-file,
code-centric} analysis. LSCPS differs in three concrete ways:

\begin{enumerate}[topsep=2pt, itemsep=2pt]
\item \textbf{Multi-file capability aggregation.}
  LSCPS extracts $\mathcal{C}_{\text{act}}$ across SKILL.md manifests,
  Python helpers, YAML configs, and JSON schemas jointly (up to 17.2 files
  per skill on average). Bandit analyzes files independently and misses
  exfiltration chains that span manifest + helper code.
\item \textbf{Dual-side extraction.}
  LSCPS extracts $\mathcal{C}_{\text{vis}}$ from \emph{natural language
  descriptions} as well as $\mathcal{C}_{\text{act}}$ from code.
  No general linter performs description-vs-code inconsistency comparison.
\item \textbf{Marketplace-specific patterns.}
  LSCPS covers 14 families including MCP permission declarations,
  autonomous-mode signals, and SKILL.md structural cues outside general linters'
  default rule sets. This yields $\approx$35\% added recall over Bandit
  (high-severity) on our T6 evaluation corpus.
\end{enumerate}

Beyond these differentiators, static extraction provides three deployment advantages:
\begin{itemize}
    \item \textbf{Evasion resistance:} Static extraction evaluates purely
      syntactically, unlike LLM judges which are vulnerable to adversarial
      prompt manipulation injected into code comments.
    \item \textbf{Deployability:} The regex-based engine executes
      marketplace-scale scanning across tens of thousands
      of skills in seconds with zero runtime overhead.
    \item \textbf{High baseline recall:} Our 300+ pattern bank identifies
      72.5\% of real-world T6 malicious pairs at the pair level,
      compared to the platform scanner's 0\%.
\end{itemize}

\paragraph{Capability weighting.}
For the static baseline (SCS), we assign heuristic
severity weights to each LSCPS family as an initial
Capability Exposure Score (CES): shell execution (3.0),
credential transmit (3.0), credential read (2.5),
network write (2.0), autonomous (1.5), file read (0.5).
In \S\ref{sec:riskweights}, we replace these heuristics
with \emph{data-driven} Risk-Weighted Mismatch (RWM) weights
learned from T6 ground truth. The learned weights diverge
significantly from the heuristic assignments
(\eg shell execution drops from 3.0 to 0.10 because it is
ubiquitous across both malicious and benign skills),
demonstrating the necessity of empirical calibration.

\subsection{Datasets}

\paragraph{\clawhub{}.}
We archived the full \clawhub{} public skill catalog in March 2026,
capturing \NskillsCH{} skills and complete version histories.
Of these, \textbf{\NskillsDMR{} skills} have locally archived
\texttt{SKILL.md} text, covering \textbf{\Ntransitions{} version transition pairs}
(on average 3.39 versions per skill).
The 705 remaining skills were deleted from the platform before
local archival completed.
\textit{Deleted-skill bias note:} Because deleted skills may disproportionately
represent platform-removed (and thus potentially higher-risk) content,
excluding them conservatively biases our IC-1 prevalence \emph{downward}---
our reported figures therefore underestimate the true inconsistency rate.

\paragraph{\skillssh{}.}
We archived \textbf{116,208 skills} from \skillssh{}.
Unlike \clawhub{}, \skillssh{} does not expose
multi-version histories in its public API, so our analysis is static-only.
We apply the identical IC framework pipeline
(\S\ref{sec:static}) to validate cross-platform generalizability.

\paragraph{Ethical considerations.}
All data is from public APIs at rate $\leq$ 0.5 req/s.
We did not execute any skill, access non-public content,
or interfere with platform operations.
Our crawler identified itself with an academic research user agent.
T6 case identifiers are withheld pending coordinated disclosure.
